\documentclass[../../main.tex]{subfiles}
\usepackage{scrextend}
\begin{document}

{\bf ［解］}

題意の円を$C: (x - 2a)^2 + (y - a)^2 = 5a^2 - 20a + 25$ とおく．

まず前半部分，円$C$が二つの定点を通ることを示す．円の方程式を$a$について整理すると
\begin{align}
 (-4x - 2y + 20)a + (x^2 + y^2 - 25) = 0
\end{align}
だから，これが $a$ についての恒等式の時，つまり
\begin{align}
\begin{cases}
-4x - 2y + 20 = 0 \\
x^2 + y^2 = 25
\end{cases}
\iff (x,y) = (3,4), (5,0) 
\end{align}
の時に与式は $a$ によらず成り立つから，もとめる $2$ 定点はこれである．

次に後半部分を考える．二つ目の円を$D: x^2+y^2=25$とする．
二つの円$C, D$の中心間距離 $d$ は 
\begin{align}
 d = \sqrt{(2a)^2 + a^2} = \sqrt{5}|a| 
\end{align}
である．
二つの円が接するパターンは外接と内接があるから場合分けして考える．

% \begin{addmargin}[2em]{0pt}
\subparagraph{$1^\circ$ 外接する時}
二つの円の半径の和が中心間距離$d$に等しいから
\begin{align}
& \sqrt{5} + \sqrt{5}\sqrt{a^2 - 4a + 5} = \sqrt{5}|a| \\
& \sqrt{a^2 - 4a + 5} = |a| - 1
\end{align}
である．$|a| \ge 1$ のもとで両辺二乗して整理すると $a = 2$ を得る．

\subparagraph{$2^\circ$ 内接する時}
二つの円の半径の差が中心間距離$d$に等しいから
\begin{align}
| \sqrt{5} - \sqrt{5}\sqrt{a^2 - 4a + 5} | = \sqrt{5}|a| 
\end{align}
両辺 $0$ 以上だから二乗して
\begin{align}
a^2 - 4a + 5 + 1 - 2\sqrt{a^2 - 4a + 5} = a^2 \\
\sqrt{a^2 - 4a + 5} = 3 - 2a
\end{align}
$a \le \dfrac{3}{2}$ のもとで両辺二乗して
\begin{align}
   & a^2 - 4a + 5 = 4a^2 - 12a + 9 \\
   & 3a^2 - 8a + 4 = 0 \\
   & (3a - 2)(a - 2) = 0 \\
\therefore 
   & a = \dfrac{2}{3}, 2
\end{align}
$a \le \dfrac{3}{2}$ をみたすのは $a = \dfrac{2}{3}$である．
% \end{addmargin}
\casesend

以上から求める値は
\begin{align}
 a = \dfrac{2}{3}, 2 
\end{align}
である．

\begin{figure}[htb]
\begin{center}
  \begin{tikzpicture}[scale=0.8, >=stealth]
  % 軸の描画
  \draw[->, thick] (-6, 0) -- (8, 0) node[right] {$x$};
  \draw[->, thick] (0, -6) -- (0, 7) node[above] {$y$};
  \node[below left] at (0,0) {$O$};

  % 円D
  \def\rD{2.236} % sqrt(5)
  \coordinate (CenterD) at (0,0);

  % 円C (a = 2/3)
  % \coordinate (CenterC) at (1.333, 0.667); % (4/3, 2/3)
  % \def\rC{5.218} % sqrt(5)*sqrt((2/3)^2 - 4*(2/3) + 5) = 7*sqrt(5)/3 ≒ 5.218

  \coordinate (CenterC) at (4, 2); 
  \def\rC{2.236} % sqrt(5)


  % 2つの定点 (3,4) と (5,0)
  \coordinate (P1) at (3,4);
  \coordinate (P2) at (5,0);

  % 円 D (x^2 + y^2 = 25)
  \draw[thick, blue!80!black] (CenterD) circle (\rD);
  \node[blue!80!black, font=\small] at (-0.8, -1.2) {$D$};
  \fill[blue!80!black] (CenterD) circle (2pt) node[below left] {$O$};

  % 円 C (a = 2/3)
  \draw[thick, red!80!black] (CenterC) circle (\rC);
  \node[red!80!black, font=\small, right] at (4.8, 5.2) {$C\;(a=2)$};
  \fill[red!80!black] (CenterC) circle (2pt) node[below right, font=\small] {$\left(4,2\right)$};

  % 中心間を結ぶ線分 d
  \draw[thick, purple, dashed] (CenterD) -- (CenterC);


  % 定点 (3,4), (5,0) のプロットとラベル
  \fill[black] (P1) circle (2.5pt) node[above right] {$(3,4)$};
  \fill[black] (P2) circle (2.5pt) node[below right] {$(5,0)$};

\end{tikzpicture}
\caption{$a=2$で二つの円が外接する場合}
\end{center}
\end{figure}

\begin{figure}[htb]
\begin{center}
  \begin{tikzpicture}[scale=0.8, >=stealth]
  % 軸の描画
  \draw[->, thick] (-6, 0) -- (8, 0) node[right] {$x$};
  \draw[->, thick] (0, -6) -- (0, 7) node[above] {$y$};
  \node[below left] at (0,0) {$O$};

  % 円D
  \def\rD{2.236} % sqrt(5)
  \coordinate (CenterD) at (0,0);

  % 円C (a = 2/3)
  \coordinate (CenterC) at ({4/3}, {2/3});
  \def\rC{sqrt(125)/3} % sqrt(20/9-40/3+25)=sqrt(125/9)


  % 2つの定点 (3,4) と (5,0)
  \coordinate (P1) at (3,4);
  \coordinate (P2) at (5,0);

  % 円 D (x^2 + y^2 = 25)
  \draw[thick, blue!80!black] (CenterD) circle (\rD);
  \node[blue!80!black, font=\small] at (-0.8, -1.2) {$D$};
  \fill[blue!80!black] (CenterD) circle (2pt) node[below left] {$O$};

  % 円 C (a = 2/3)
  \draw[thick, red!80!black] (CenterC) circle (\rC);
  \node[red!80!black, font=\small, right] at (4.8, 5.2) {$C\;(a=2/3)$};
  \fill[red!80!black] (CenterC) circle (2pt) node[below right, font=\small] {$\left(\dfrac{4}{3},\dfrac{2}{3}\right)$};

  % 中心間を結ぶ線分 d
  \draw[thick, purple, dashed] (CenterD) -- (CenterC);


  % 定点 (3,4), (5,0) のプロットとラベル
  \fill[black] (P1) circle (2.5pt) node[above right] {$(3,4)$};
  \fill[black] (P2) circle (2.5pt) node[below right] {$(5,0)$};

\end{tikzpicture}
\caption{$a=\dfrac{2}{3}$で二つの円が内接する場合}
\end{center}
\end{figure}



\newpage

\end{document}
